
\section{Organization of the work}
The second chapter serves as a more formal and detailed introduction necessary for further considerations. In Section \ref{sec:graph_search_problem} we recall the basic notions of graph theory required for the analysis and in Section \ref{sec:optimization_problems} we define optimization problems as well as the approximation algorithms and their classes. Then, in Section \ref{sec:graph_search_problem} we formally state the general Graph Search Problem we are interested in throughout this work.

In the third chapter we focus ourselves on the formal analysis of the worst-case variants of the search problem on trees, including the presentation of the most interesting algorithmic results known for this problem as well as few new contributions. We begin with the exact dynamic programming algorithm for the uniform cost query  model due to \cite{OnakParys2006GenOfBSSInTsAndFLikePosets,Mozes_Onak2008FindOptTSStartInLinTime} in Section \ref{sec:trees_uniform_costs}.
In Section \ref{sec:log_n_log_log_n_approx} we show an $O\br{\log n/\log \log n}$-approximation algorithm by \cite{Cicalese2016OnTSPwNonUniCost} adapted to the vertex query model which is based on an exact exponential time procedure running in $O\br{2^nn}$ time. Then, in Section \ref{sec:sqrt_log_n_approx} we show an $O\br{\sqrt{\log n}}$-approximation algorithm following a similar idea by \cite{dereniowski2017ApproxSsForGeneralBSinWTs}. This result is based on the QPTAS presented in the later work. However, the construction of the original procedure is counter-intuitive. To aleviate this inconvienience, in Section \ref{qptasSection} we restate the QPTAS by showing a novel dynamic program procedure based on the algorithm used for $T||V||\sum C_i$ \cite{Jacobs2010OnTheComplexSearchInTsAvg,Angelidakis2018ShortestPQ}. However, it should be noted that this is not a straightforward adaptation and required significant strenghtening of the basic procedure. As a second corollary of the aforementioned QPTAS, in Section \ref{sec:log_log_n_approx} we describe how to use it to obtain an $O\br{\log \log n}$-approximation algorithm for the Tree Searh Problem running in time $k^{O\br{\log k}}\cdot n^{O\br{1}}$, where $k$ is the $k$-up modularity of the cost function.

In the fourth chapter we present the results of the computational experiments conducted to evaluate the practical performance of the algorithms presented in previous chapter. In Section~\ref{sec:experiments_overview} we describe the experimental setup, including: tree generation methods and cost distributions used. For the uniform cost model, in Section \ref{sec:trees_uniform_costs} we evaluate the performance of the exact algorithm on random trees, bounded degree trees, and bounded diameter trees. For the non-uniform cost model, in Sections \ref{sec:random_c_uniform} to \ref{sec:bounded_diameter_c_exponential} we measure the approximation quality and runtime of the $O\br{\log n/\log \log n}$-approximation, $O\br{\sqrt{\log n}}$-approximation, and $O\br{\log \log n}$-approximation algorithms across various tree structures and cost distributions (uniform, binomial, and exponential). Additionally, we compare these approximation algorithms against the optimal solutions computed by the exact exponential time algorithm for small instances. Finally, in Section \ref{sec:exact_dp_runtime_analysis} we present regression analysis confirming that the runtime of the exact dynamic programming algorithm in fact follows an exponential function with base much lower then $2$.

The fifth chapter serves as a summary of our considerations and points the further research directions regarding this field.
